\documentclass[a4paper,12pt]{book}

% PACKAGES DE BASE
\usepackage[english,french]{babel}
\usepackage[utf8]{inputenc}
\usepackage[T1]{fontenc}

% PACKAGES ADDITIONNELS
\usepackage{color}  % pour utiliser des couleurs
\usepackage{enumitem} % pour personnaliser les puces
\usepackage{fancyhdr}  % Personnalisation en-tête et pied de page
\usepackage[top=2.5cm, bottom=2.5cm, left=2.8cm, right=2.8cm]{geometry} % Pour redimensionner les marges
\usepackage{graphicx}  % Pour l'insertion d'images
\usepackage[hyphens]{url}  % Pour insérer des url. l'option 'hyphens' permet de ne pas afficher de barre oblique au niveau de la césure d'une url.
\usepackage{listings}  % Pour définir un environnement de citation de code
\usepackage{tcolorbox}  % Pour le paramètrages des box
\usepackage{titlesec}  % Pour personnaliser les titres des chapitres
\usepackage{verbatim}  % pour les citations de code
\usepackage{xcolor}  % pour utiliser des couleurs
% Package à supprimer
\usepackage{lipsum}
% Package à placer en dernier
\usepackage{hyperref}

% DEFINITION DES COULEURS
\definecolor{backcolour}{rgb}{0.95,0.95,0.92}
\definecolor{codegreen}{rgb}{0,0.6,0}
\definecolor{pythonblue}{RGB}{54,117,171}
\definecolor{pythonyellow}{RGB}{220,176,30}
\definecolor{graylight}{RGB}{231,231,231}

% PERSONNALISER LES TITRES DES CHAPITRES
\titleformat{\chapter}[frame]
  {\normalfont\Huge\bfseries\color{pythonblue}}
  {\textcolor{pythonyellow}{Challenge N°}\ \textcolor{pythonyellow}{\thechapter}}
  {25pt}
  {\Huge\centering}
  

% PERSONNALISER LES TITRES DES SECTIONS
\titleformat{\section}
  {\normalfont\LARGE\bfseries\color{pythonblue}}
  {\textcolor{pythonyellow}{\thesection}}{1em}{}
  {\LARGE} 
 
% CREATION DE COMMANDES
\newcommand{\bigger}{\fontsize{80}{100}\selectfont}  % Police de grande taille

% ENVIRONNEMENT POUR DU CODE PYTHON:
\lstset{
    language=Python,
    aboveskip=5mm,  % Marge avant l'insertion du code
    belowskip=5mm,  % Marge après l'insertion du code
    backgroundcolor=\color{backcolour},
    keywordstyle=\color{red}, % Couleur mots clés
    commentstyle=\color{olive}, % Coleur commentaires  
    stringstyle=\color{codegreen},  % Style des strings
    basicstyle=\small, % Taille police
    basicstyle=\ttfamily, % Style police
    numbers=left, % Numéros de ligne à gauche
    numberstyle=\scriptsize, % Taille numéros de ligne 
    numbersep=10pt, % Espace entre numéro de ligne et code
    showspaces=false,  % Voir les espaces             
    showstringspaces=false,  % Voir les espaces des strings
    showtabs=false,  % Voir les tabulations               
    tabsize=4,  % Tabulations
    frame=tb,  % Cadre en haut et en bas
    rulecolor=\color{gray},  % Couleur du cadre 
    captionpos=b,  % Position de la légende du code
    literate= % Pour une gestion de l'utf8
        {â}{{\^a}}1 {à}{{\`a}}1 {ç}{{\c c}}1 {é}{{\'e}}1
        {è}{{\`e}}1 {ê}{{\^e}}1 {î}{{\^i}}1 {ô}{{\^o}}1
        {É}{{\`E}}1
        {Ò}{{\`O}}1
}

%%% TEXTE %%%

\begin{document}

% PAGE DE COUVERTURE
\thispagestyle{empty} % Supprime la numérotation de la première page
\begin{center}
\begin{tcolorbox}[colback=graylight, boxrule=2pt, sharp corners, toprule=2pt, bottomrule=2pt] %[frame=double, colback=#f2f2f2]
\begin{minipage}{\textwidth}
    \centering
    \vspace{1.8cm}
    \bigger \textbf{\textcolor{pythonblue}{Challenges Python}} \\
    \vspace{1cm}
    \Large \textbf{\textcolor{pythonyellow}{(Discord \og Docstring\fg{})}}
    
    \vspace{3cm}
    \includegraphics[width=13cm]{IMG/logo_python.jpeg} 
    
    \vspace{6cm}
    
    \Large \textbf{\textcolor{pythonblue}{@bucdany}} \\
    \vspace{2cm}
    2023
    \vspace{1.2cm}
\end{minipage}
\end{tcolorbox}
\end{center}

% PREMIERE  (vide - première page numérotée)
\clearpage 
\setcounter{page}{1}
\pagestyle{plain}

% TABLE DES MATIERES
\tableofcontents

% DEBUT TEXTE
\chapter*{Avant propos}
\lipsum[1-3]

% CHAPITRE 1
\chapter{Compter le nombre de voyelles}
\vspace{2cm}
\section{Énoncé}
\medskip

Ce premier challenge est très simple, il est de niveau \og débutant\fg{}, mais si vous avez plus d'expérience, vous pouvez essayer de trouver de belles astuces pour un code propre, rapide et concis.
\medskip

Ici, il va nous falloir créer une fonction \verb|nb_voyelles(phrase: str)->int| qui retourne le résultat du nombre total de voyelles dans une phrase passée en paramètre.
\medskip

\subsection*{Conditions}
\begin{itemize}
	\item[-] Les voyelles sont : \texttt{aeiou}, \texttt{y} n’est pas pris en compte.
	\item[-] Les voyelles accentuées ne sont pas prises en compte.
	\item[-] La phrase passée en paramètre doit être écrite en minuscule.
	\item[-] Une chaîne vide, passée en paramètre, doit renvoyer \texttt{0}.
\end{itemize}
\medskip

\subsection*{Exemples}
\begin{itemize}
	\item[\textbullet] \verb|nb_voyelles("bonjour, comment allez-vous ?")| doit retourner \texttt{9}.
	\item[\textbullet] \verb|nb_voyelles("je vais à paris")| doit retourner \texttt{5}.
	\item[\textbullet] \verb|nb_voyelles("docstring")| doit retournet \texttt{2}.
	\item[\textbullet] \verb|nb_voyelles("")| doit retourner \texttt{0}.
\end{itemize}
\medskip

\end{document}